\section{Availability and Limitations}

The artifact repository contains the \tool Python package, adapter scripts,
metric runners, generated CSV, Markdown, and \LaTeX{} tables, provenance logs,
the BibTeX/source audit used for this draft, an MIT license, and a reviewer
quickstart. Large local artifacts are kept out of git and described through
manifests so that reviewers can distinguish public reproducibility assets from
local caches.

\begin{table}[t]
\caption{Reviewer-facing artifact package. The table maps repository assets to
verification actions and claims rather than asking reviewers to reproduce all
upstream tokenizer training from scratch.}
\label{tab:artifact-package}
\small
\setlength{\tabcolsep}{3pt}
\begin{tabularx}{\columnwidth}{@{}p{0.42\columnwidth}Y@{}}
\toprule
Reviewer action & Checked claim \\
\midrule
Read quickstart and license & MIT-licensed resource with local verification
commands. \\
Run tests and table builders & D1--D5a metrics and optional D6 paths execute
from normalized artifacts. \\
Compare generated CSVs with Tables~\ref{tab:method-coverage}
and~\ref{tab:musical-diagnostic} & Main numeric claims are traceable to
generated files. \\
Inspect manifest and claim audit & Rows are labeled as runnable, controlled,
optional, or literature-only evidence. \\
Open sanity, scale, churn, and MovieLens CSVs & Metric sensitivity, larger-item
handling, D6 churn, and non-Amazon portability are documented. \\
\bottomrule
\end{tabularx}
\end{table}

The current evidence supports a conservative resource-demo claim. It does not
support several stronger claims. \tool does not reproduce TIGER end-to-end; the
GRID Musical row is a controlled feature-text export. The ReSID Sports balanced
GAOQ path was stopped because constrained k-means was CPU-bound under the
available window, so the main ReSID evidence is the bounded Musical export.
CARD is repaired only as a fallback/stressor path and should not be presented
as faithful CARD reproduction~\cite{wei2026card}. D2 is a collision profile,
not strict causal harm. D3 is a diagnostic proxy and is not yet proven
monotonic with Recall@K or NDCG. D5a is a structural cost proxy without real
generator-output latency, and D6 remains optional drift evidence.

These boundaries are intentional. The paper contributes an audit surface for
tokenizer artifacts, not a new \sid tokenizer, a complete
generative-recommender evaluation, or an industrial online-impact study.
Future versions should add interaction-qualified collision harm, generator
output diagnostics, stronger same-dataset method coverage, and unified search
and recommendation artifact checks~\cite{penha2025jointsid}. Industrial SID
deployment studies further motivate tracking these artifact properties as
first-class engineering objects~\cite{ju2026snapchatsid}.

The most important limitation is interpretive rather than computational.
Artifact diagnostics can reveal pressure points before expensive generator
training, but they do not replace ranking evaluation. A tokenizer with cleaner
capacity allocation can still be hard for a sequence model to predict, and a
tokenizer with collisions can sometimes be rescued by downstream ranking or
reranking components. For that reason, \tool is best used as a preflight and
debugging layer: it tells method authors what kind of artifact they produced,
which slices deserve downstream evaluation, and which claims are unsafe to make
from Recall@K alone.

\paragraph{Reviewer workflow.}
The intended reviewer workflow is short. First, inspect the manifest to decide
which rows are named-method evidence, controlled stressors, optional
extensions, or literature-only context. Second, rerun the metric command on the
normalized artifacts and compare the generated CSV and \LaTeX{} tables with
the paper-facing tables. Third, open the failure slices for any suspicious row:
items sharing a full code, prefixes with unusually high fan-out, or tail items
with low unique-code ratios. This workflow is deliberately lighter than
retraining every tokenizer. It tests whether the resource is useful as an audit
layer for artifacts that already exist.

\paragraph{Claim discipline.}
The resource separates three levels of statements. A \emph{mapping statement}
is directly supported by the SID-assignment table: for example, the number of
unique full codes or the full-collision rate in
Table~\ref{tab:musical-diagnostic}. A \emph{diagnostic statement} combines the
mapping with metadata or interaction slices: for example, D3 collaborative
prefix recovery or D4 tail allocation. A \emph{system-quality statement}
requires downstream generator training, ranking evaluation, or serving traces.
The current paper makes the first two kinds of statements and explicitly does
not make the third. This separation is the main design rule that keeps the
resource useful even when upstream tokenizer releases differ in how much code,
checkpoint state, or generator output they expose.

\paragraph{Why this is still resource-worthy.}
A resource paper does not need to prove that one tokenizer is better than
another. It needs to make a reusable artifact easier to inspect, compare, and
extend. \tool contributes that layer for SID tokenizers: adapters translate
heterogeneous exports into one contract, validators prevent unjoinable mappings
from entering the evidence table, diagnostics expose failure modes that are
orthogonal to final Recall@K, and generated tables make the audit repeatable.
The same design also gives future tokenizer papers a checklist for what to
release: item-to-\sid mappings, item metadata, interaction slices when
possible, generator outputs when claiming decoding behavior, and enough
provenance to distinguish a faithful reproduction from a controlled proxy.

\paragraph{Maintenance plan.}
The public artifact is organized so that new tokenizer adapters can be added
without changing the diagnostic definitions. Each adapter is expected to emit
the same normalized tables, declare its evidence role, and record the source of
item features or checkpoints. New diagnostics should be added as separate D
modules only when they require new input contracts; otherwise they should
extend the existing CSV schema. This keeps the resource useful for future SID
papers without turning the main package into a collection of one-off scripts.
Versioned manifests and fixed-name ``latest'' reports preserve both stable
entry points and timestamped provenance for reviewers.
